% GENERATED FILE. Edit canonical lessons, then run npm run generate:books.
% canonical-source-hash: fnv1a64:94779d32984dade2
% canonical-lessons: ES-C544-nacimiento, ES-C544-parto, ES-C544-dar-a-luz, ES-C544-embarazo, ES-C544-recien-nacido

\chapter{Un bebé en camino}
\label{ch:es-un-bebe-en-camino}
\hlchaptermodality{544}

\begin{chapteropening}
\textbf{By the end of this chapter:} \emph{I can talk about a pregnancy and a birth: the newborn baby and how the birth went.}

Complete the last lesson of chapter 544: I can talk about a pregnancy and a birth: the newborn baby and how the birth went.
\end{chapteropening}

\section[el nacimiento]{el nacimiento — the birth}

\label{lesson:ES-C544-nacimiento}



\begin{quote}
Before the new one: say the Spanish for to bring forward, then the Spanish for to remain.
\end{quote}

\subsection*{You'll want to know: el nacimiento}
\textbf{el nacimiento} — \textquotedblleft{}the birth\textquotedblright{}.

\begin{quote}
Celebramos el nacimiento de mi sobrina.
\end{quote}

\begin{quote}
\emph{We celebrated my niece's birth.}
\end{quote}

\emph{La fecha de nacimiento} — the date of birth.

\begin{cousinweb}
\emph{nacimiento} is built on \emph{nacer} 'to be born' with the ending \emph{-miento}. \emph{Nacer} comes from Latin \emph{nāscī} 'to be born' (Lewis \& Short).
\end{cousinweb}

\begin{grammarlens}[title={What you've built}]
One more word for telling, explaining and describing at B1.
\end{grammarlens}

\subsection*{Your turn}
\begin{itemize}
\raggedright
  \item \emph{Say it:} \emph{el nacimiento}
  \item \emph{Say it:} \emph{el nacimiento}, once more
  \item \emph{Say it:} a sentence of your own with \emph{el nacimiento}
  \item \emph{Recall:} say the Spanish for to bring forward, then the Spanish for to remain, then say \emph{el nacimiento} again
\end{itemize}

\begin{tcolorbox}[breakable,colback=teal!4,colframe=teal!35!black,title={Before you move on}]
What does el nacimiento mean? (\textquotedblleft{}The birth\textquotedblright{}.) Say it once more.
\end{tcolorbox}
\section[el parto]{el parto — the delivery (of a baby), the birth}

\label{lesson:ES-C544-parto}



\begin{quote}
Before the new one: say the Spanish for to remain, then the Spanish for the birth.
\end{quote}

\subsection*{You'll want to know: el parto}
\textbf{el parto} — \textquotedblleft{}the delivery (of a baby), the birth\textquotedblright{}.

\begin{quote}
El parto fue largo pero todo salió bien.
\end{quote}

\begin{quote}
\emph{The birth was long but everything went well.}
\end{quote}

\emph{Estar de parto} — to be in labour.

\begin{cousinweb}
\emph{parto} is inherited from Latin \emph{partus} 'a birth' (EtymDB), built on \emph{parere} 'to bring forth, to give birth'. It is not the \emph{parere} 'to appear' behind \emph{parecer}.
\end{cousinweb}

\begin{grammarlens}[title={What you've built}]
One more word for telling, explaining and describing at B1.
\end{grammarlens}

\subsection*{Your turn}
\begin{itemize}
\raggedright
  \item \emph{Say it:} \emph{el parto}
  \item \emph{Say it:} \emph{el parto}, once more
  \item \emph{Say it:} a sentence of your own with \emph{el parto}
  \item \emph{Recall:} say the Spanish for to remain, then the Spanish for the birth, then say \emph{el parto} again
\end{itemize}

\begin{tcolorbox}[breakable,colback=teal!4,colframe=teal!35!black,title={Before you move on}]
What does el parto mean? (\textquotedblleft{}The delivery (of a baby), the birth\textquotedblright{}.) Say it once more.
\end{tcolorbox}
\section[dar a luz]{dar a luz — to give birth}

\label{lesson:ES-C544-dar-a-luz}



\begin{quote}
Before the new one: say the Spanish for the birth, then the Spanish for the delivery.
\end{quote}

\subsection*{You'll want to know: dar a luz}
\textbf{dar a luz} — \textquotedblleft{}to give birth\textquotedblright{}.

\begin{quote}
Dio a luz a una niña el martes.
\end{quote}

\begin{quote}
\emph{She gave birth to a girl on Tuesday.}
\end{quote}

\emph{Dar a luz a} + baby. Word for word, to give to the light.

\begin{cousinweb}
\emph{dar a luz} is built on \emph{dar}, inherited from Latin \emph{dare} 'to give', and \emph{luz}, inherited from Latin \emph{lūx} 'light' (Lewis \& Short). To give birth is to give the child to the light.
\end{cousinweb}

\begin{grammarlens}[title={What you've built}]
One more word for telling, explaining and describing at B1.
\end{grammarlens}

\subsection*{Your turn}
\begin{itemize}
\raggedright
  \item \emph{Say it:} \emph{dar a luz}
  \item \emph{Say it:} \emph{dar a luz}, once more
  \item \emph{Say it:} a sentence of your own with \emph{dar a luz}
  \item \emph{Recall:} say the Spanish for the birth, then the Spanish for the delivery, then say \emph{dar a luz} again
\end{itemize}

\begin{tcolorbox}[breakable,colback=teal!4,colframe=teal!35!black,title={Before you move on}]
What does dar a luz mean? (\textquotedblleft{}To give birth\textquotedblright{}.) Say it once more.
\end{tcolorbox}
\section[el embarazo]{el embarazo — the pregnancy}

\label{lesson:ES-C544-embarazo}



\begin{quote}
Before the new one: say the Spanish for the delivery, then the Spanish for to give birth.
\end{quote}

\subsection*{You'll want to know: el embarazo}
\textbf{el embarazo} — \textquotedblleft{}the pregnancy\textquotedblright{}.

\begin{quote}
Tuvo un embarazo tranquilo.
\end{quote}

\begin{quote}
\emph{She had an easy pregnancy.}
\end{quote}

It once meant an obstacle, and in formal Spanish it still can.

\begin{cousinweb}
\emph{embarazo} is built on \emph{embarazar}, first 'to get in the way, to hinder'. The origin of \emph{embarazar} is not settled; it is often linked to Portuguese \emph{embaraçar} and an old word \emph{baraça} 'a cord'.
\end{cousinweb}

\begin{grammarlens}[title={What you've built}]
One more word for telling, explaining and describing at B1.
\end{grammarlens}

\subsection*{Your turn}
\begin{itemize}
\raggedright
  \item \emph{Say it:} \emph{el embarazo}
  \item \emph{Say it:} \emph{el embarazo}, once more
  \item \emph{Say it:} a sentence of your own with \emph{el embarazo}
  \item \emph{Recall:} say the Spanish for the delivery, then the Spanish for to give birth, then say \emph{el embarazo} again
\end{itemize}

\begin{tcolorbox}[breakable,colback=teal!4,colframe=teal!35!black,title={Before you move on}]
What does el embarazo mean? (\textquotedblleft{}The pregnancy\textquotedblright{}.) Say it once more.
\end{tcolorbox}
\section[el recién nacido]{el recién nacido — the newborn baby}

\label{lesson:ES-C544-recien-nacido}



\begin{quote}
Before the new one: say the Spanish for to give birth, then the Spanish for the pregnancy.
\end{quote}

\subsection*{You'll want to know: el recién nacido}
\textbf{el recién nacido} — \textquotedblleft{}the newborn baby\textquotedblright{}.

\begin{quote}
El recién nacido pesó tres kilos.
\end{quote}

\begin{quote}
\emph{The newborn baby weighed three kilos.}
\end{quote}

It agrees: \emph{una niña recién nacida}.

\begin{cousinweb}
\emph{el recién nacido} is built on \emph{recién}, the short form of \emph{reciente}, from Latin \emph{recēns} 'fresh, new', and \emph{nacido}, the participle of \emph{nacer}, from Latin \emph{nāscī} 'to be born' (Lewis \& Short).
\end{cousinweb}

\begin{grammarlens}[title={What you've built}]
One more word for telling, explaining and describing at B1.
\end{grammarlens}

\subsection*{Your turn}
\begin{itemize}
\raggedright
  \item \emph{Say it:} \emph{el recién nacido}
  \item \emph{Say it:} \emph{el recién nacido}, once more
  \item \emph{Say it:} a sentence of your own with \emph{el recién nacido}
  \item \emph{Recall:} say the Spanish for to give birth, then the Spanish for the pregnancy, then say \emph{el recién nacido} again
\end{itemize}

\begin{tcolorbox}[breakable,colback=teal!4,colframe=teal!35!black,title={Before you move on}]
What does el recién nacido mean? (\textquotedblleft{}The newborn baby\textquotedblright{}.) Say it once more.
\end{tcolorbox}
